The discovery of novel stable materials is a cornerstone of technological progress, yet the vast search space of elemental combinations remains a major bottleneck for human-driven trial-and-error discovery. While high-throughput computational methods and graph neural networks have mapped millions of structures, intelligently traversing this space to identify targetable compositions remains challenging. In this work, we investigate whether large language models (LLMs) can autonomously discover stable elemental combinations by leveraging their intrinsic chemical reasoning and active feedback loops. We propose \ours, an autonomous framework that utilizes \gpt to hypothesize stable ternary material systems, which are then validated against the \gnome dataset. Our results show that \ours with \zeroshot prompting achieves a 50\% hit rate and proposes combinations with significantly lower formation energies (average -1.30 eV/atom) compared to random sampling (-0.82 eV/atom), representing a medium effect size (Cohen's $d = 0.58$). Furthermore, by employing an \activelearning loop, the system identifies exceptionally stable materials (average -2.33 eV/atom), such as alkaline-earth titanates. These findings demonstrate that LLM-driven autonomous agents can systematically accelerate material invention by bridging the gap between high-level reasoning and high-throughput physical databases.
